% One-time imported draft. Editable native source; no Markdown regeneration.
\section{Step 13 Observe a single transaction in
Debug}\label{005-seven-segment-counter-s13-step-13-observe-a-single-transaction-in-debug}

A fast Run can complete before every glyph is repainted. A breakpoint
gives the program a stable state while the tool refreshes. The labels
distinguish the just-written value from the updated loop counters.

\subsection{Write or
inspect}\label{005-seven-segment-counter-s13-write-or-inspect}

Set a breakpoint on the first executable line below \texttt{displayed}.
Debug the default fixture and inspect the right glyph and raw mask.
Resume to later hits, or stop below \texttt{iteration\_done} to inspect
the updated counters.

\subsection{Expected
observation}\label{005-seven-segment-counter-s13-expected-observation}

After the first store the right mask is 3f. At the eleventh
\texttt{displayed} hit it is 77 for A. At done it is 71 for F, and the
left unit remains blank.

\subsection{Check your
understanding}\label{005-seven-segment-counter-s13-check-your-understanding}

Does a completed Run showing F prove that the intermediate glyphs were
written?

\begin{itemize}
\tightlist
\item
  \textbf{A} No; the store sequence needs separate evidence.
\item
  \textbf{B} Yes; the final glyph encodes the entire history.
\item
  \textbf{C} Yes; the UI always repaints after every instruction.
\end{itemize}
